\documentclass{report}
\usepackage{amsfonts, amsmath}
\newcommand\R{{\mathbb R}}
\newcommand\Q{{\mathbb Q}}
\newcommand\N{{\mathbb N}}
\pagestyle{empty}
\begin{document}
\large
\centerline{\Huge Analisi Matematica II}
\renewcommand{\thefootnote}{ } 
\footnote{\sl Avete 2 ore di tempo. Ogni esercizio
vale nove punti. La sufficienza si ottiene con un punteggio $\geq$
18. Solo le risposte chiaramente giustificate saranno prese in
considerazione.}
\renewcommand{\thefootnote}{\arabic{footnote}} 
\vskip.1cm
\centerline{\Large Prova di Autovalutazione}
\vskip1cm
\begin{enumerate}
\item Calcolare, per $n,m\in\N$,
\[
\int^{2\pi}_0\sin2n\theta\,\cos m\theta\, d\theta.
\]

\item Si calcoli il valore di
\[
\sum_{n=1}^\infty \frac 1{n^2}
\]
con una precisione superiore a $10^{-2}$.
\item Per ogni $p\in\N$ si calcoli il raggio di convergenza della
serie 
\[
\sum_{n=0}^\infty n^p x^n.
\]

\item Si calcolino i seguenti limiti
\[
\lim_{L\to\infty}\int_2^L\frac{1}{x(\ln x)^2} dx\,,
\]\[
\lim_{\varepsilon\to 0}\int^1_\varepsilon \frac 1{\sqrt{x}} dx.
\]
\end{enumerate}
\end{document}
